\begin{solapartat}
La freqüència llindar d'un metall és la freqüència mínima que ha de tenir una radiació electromagnètica, per a què els seus fotons puguin arrencar electrons d'aquest metall, per efecte fotoelèctric. \punts{0,2}

A partir de la gràfica veiem que l'energia cinètica en funció de la freqüència ve donada per la recta:
\[ E_c = 3,72 + \frac{7,86 - 3,72}{3 - 2}\left(\frac{f}{10^{15}} - 2\right)\ \mathrm{eV}\ \text{\punts{0,3}} \]

A partir de l'expressió anterior obtindrem la freqüència llindar fent que l'energia cinètica sigui zero: \punts{0,2}
\[ 0 = 3,72 + 4,14\left(\frac{f_\textit{llindar}}{10^{15}} - 2\right)\ \text{\punts{0,1}} \Rightarrow f_\textit{llindar} = \left(-\frac{3,72}{4,14} + 2\right)\cdot 10^{15} = 1,10\cdot 10^{15}\ \mathrm{Hz}\ \text{\punts{0,2}} \]
\end{solapartat}

\begin{solapartat}
La constant de Planck la podem trobar a partir del pendent de la recta representada: \punts{0,2}
\[ h = \frac{(7,86 - 3,72)\ \mathrm{eV}}{(3 - 2)\ \mathrm{s^{-1}}}\times\frac{1,6\cdot 10^{-19}}{1\mathrm{eV}} = 6,62\cdot 10^{-34}\ \mathrm{Js}\ \text{\punts{0,2}} \]

\[ E_c = \frac{1}{2}mv^2 = \frac{hc}{\lambda} - h\,f_\textit{llindar}\ \text{\punts{0,2}} \Rightarrow v = \sqrt{\frac{2}{m}\left\{\frac{hc}{\lambda} - h\,f_\textit{llindar}\right\}}\ \text{\punts{0,2}} \Rightarrow \]
\begin{align*}
v &= \sqrt{\frac{2}{9,11\cdot 10^{-31}\mathrm{kg}}\left\{\frac{6,62\cdot 10^{-34}\mathrm{J\cdot s}\times 3\cdot 10^{8}\ \mathrm{m/s}}{1,2\cdot 10^{-7}\ \mathrm{m}} - 4,56\ \mathrm{eV}\,\frac{1,6\cdot 10^{-19}\ \mathrm{J}}{1\mathrm{eV}}\right\}}\\
&= 1,43\cdot 10^{6}\ \mathrm{m/s}\ \text{\punts{0,2}}
\end{align*}
\end{solapartat}
